% HD-F1-T28. GPT-6 Astra Pro; /root/transport; 22 November 2026 UTC.
\section[HD-F1-T28: every remaining Gaussian partition has a long row or column]{HD-F1-T28: every remaining Gaussian source has a row or column longer than seven eighths of its degree}
\label{sec:f1-gaussian-near-row}

For the actual residual Gaussian coefficient of T24--T27,
\[
 \max\{\lambda_1,\lambda'_1\}>\frac{7n}{8},\qquad
 4\le\min\{n-\lambda_1,n-\lambda'_1\}<\frac n8.
 \tag{F1T28.1}
\]
Consequently the original coefficient satisfies
\[
                         n\ge33,\qquad m\ge4928.
 \tag{F1T28.2}
\]
The first inequality is a uniform restriction on all partitions;
it does not enumerate one more fixed tail size.
All representations, inertia entries and geometric corner maps
are the original ones from T24--T27.

\paragraph{1. The exact retained small-cycle and odd-square budget.}
For the actual rational Specht module set
\[
 p_1=\sum_{(i,j)\in\lambda}(j-i),\qquad
 p_2=\sum_{(i,j)\in\lambda}(j-i)^2,\qquad
 N=\binom n2.
\]
The exact normalized characters of a three-cycle and a double
transposition are, with the T25 conventions,
\[
 r_3=\frac{3(p_2-N)}{n(n-1)(n-2)},\qquad
 r_{22}=\frac{4(p_1^2-3p_2+n(n-1))}
                    {n(n-1)(n-2)(n-3)}.
 \tag{F1T28.3}
\]
These are identities on the given module, obtained in T25 from
the class-sum action and its square. The actual transposition
commutator map there gives, for every nonidentity original
inertia image, the lower bound
\[
 d_\lambda=
 \max\!\left\{5,\left\lceil
 m\min\left\{\frac{1-r_3}{3},\frac{1-r_{22}}4\right\}
                           \right\rceil\right\}
 \quad\hbox{on its moving rank}.
 \tag{F1T28.4}
\]
The same applies to the two actual odd-phase inertia squares,
which T25 proves are nonidentity. Their ranks are exactly
$2a(g)-m$. Thus the two odd entries contribute at least
$m+d_\lambda$, while at least ten other original entries
contribute at least $10d_\lambda$. The original identity
$\sum a(g)=2m+6$ yields
\[
 d_\lambda\le\frac{m+6}{11}.
 \tag{F1T28.5}
\]
In particular at least one of the following necessary inequalities
holds:
\[
 r_3\ge\frac{8m-18}{11m}
 \quad\hbox{or}\quad
 r_{22}\ge\frac{7m-24}{11m}.
 \tag{F1T28.6}
\]
This inference keeps both the odd-square morphism and the
original scalar character; it does not discard scalar inertia.

\paragraph{2. Exact diagram maps bound both content moments.}
Suppose that both the first row and first column have length
at most $7n/8$. Let $r$ be the Durfee rank and write the
Frobenius arm and leg lengths as
$a_i=\lambda_i-i$, $b_i=\lambda'_i-i$ for $1\le i\le r$.
The actual bijection from the boxes of the diagram to its
diagonal hooks assigns a box $(i,j)$ with $i<j$ to the
arm of diagonal box $(i,i)$ with coordinate $j-i$, and
a box with $j<i$ to the leg of $(j,j)$ with coordinate
$i-j$. It gives exactly
\[
\begin{split}
 n&=r+\sum_i(a_i+b_i),\\
 p_2&=\sum_i\bigl(F(a_i)+F(b_i)\bigr),
 \qquad F(t)=t(t+1)(2t+1)/6.
\end{split}
 \tag{F1T28.7}
\]
All arm and leg lengths lie in $[0,U]$, where
$U=7n/8-1$, and their sum is at most $n-1$.
The function $F$ is increasing and convex on this interval.
For nonnegative numbers with fixed sum, transferring mass
from a smaller coordinate to a larger coordinate, until
one becomes zero or the larger becomes $U$, cannot decrease
the sum of their $F$ values: its derivative is
$F'(y+t)-F'(x-t)\ge0$ for $y\ge x$.
Repeat this operation. Since $n-1\le2U$ for $n\ge22$,
increasing the total mass to $n-1$ if necessary and then
concentrating it proves
\[
          p_2\le F(7n/8-1)+F(n/8).
 \tag{F1T28.8}
\]
This is an inequality on the exact hook coordinates of
the existing diagram; no other representation is substituted.

For the first moment put $M=\lfloor7n/8\rfloor$.
Every partition with first row at most $M$ has content
sum at most that of $(M,n-M)$, as follows by an explicit
sequence of box moves. Move a last box of the last nonempty
lower row to the first row until that row has length $M$.
Every move preserves a partition and strictly increases
the content sum. Then move last boxes of rows below the
second row to the second row until all $n-M$ remaining
boxes are there. These moves again increase the content
sum, and $n-M\le M$ ensures the first-row constraint.
The same argument applied to the actual transposed diagram,
whose contents are the negatives of the original ones, gives
the bound in both signs. Hence
\[
 |p_1|\le N-(n-M)(M+1)
       \le N-\frac n8\left(\frac{7n}8+1\right)
       =\frac{n(25n-40)}{64}.
 \tag{F1T28.9}
\]
For the second inequality the function
$N-(n-x)(x+1)$ is increasing for $x\ge(n-1)/2$,
and $M\ge n/2$ here.
Finally the Cauchy--Schwarz inequality on the actual
$n$ box contents gives
\[
                              p_2\ge p_1^2/n.
 \tag{F1T28.10}
\]

\paragraph{3. Both necessary character inequalities fail.}
Substitution of (F1T28.8) into the first formula of
(F1T28.3) gives
\[
 r_3\le U_3(n):=
       \frac{43n^2-168n+128}{64(n-1)(n-2)}.
 \tag{F1T28.11}
\]
For the second formula, first use (F1T28.10) and then
(F1T28.9). The coefficient of $p_1^2$ is positive
for $n>3$, so these substitutions have the indicated
directions. They give
\[
\begin{split}
 r_{22}
 &\le\frac{4p_1^2}{n^2(n-1)(n-2)}
                 +\frac4{(n-2)(n-3)}\\
 &\le U_{22}(n):=
 \frac{625n^3-3875n^2+11696n-8896}
      {1024(n-1)(n-2)(n-3)} .
\end{split}
 \tag{F1T28.12}
\]
By T27, $m\ge D_3(n):=n(n-1)(n-5)/6$.
The two right sides of (F1T28.6) are increasing functions
of $m$. The exact differences at this lower bound are
\[
\begin{split}
 \frac8{11}-\frac{18}{11D_3(n)}-U_3(n)
 &=\frac{3P_4(n)}
         {704n(n-5)(n-2)(n-1)},\\
 \frac7{11}-\frac{24}{11D_3(n)}-U_{22}(n)
 &=\frac{P_5(n)}
         {11264n(n-5)(n-3)(n-2)(n-1)},
\end{split}
 \tag{F1T28.13}
\]
where
\[
\begin{split}
P_4(n)&=13n^4+39n^3-648n^2-1664n+4608,\\
P_5(n)&=293n^5-1848n^4-47893n^3
                     +156432n^2+463040n-884736.
\end{split}
\]
Their strict positivity for every $n\ge22$ is transparent
after setting $x=n-22\ge0$:
\[
\begin{split}
P_4(22+x)
 &=13x^4+1183x^3+39678x^2+580148x+3114968,\\
P_5(22+x)
 &=293x^5+30382x^4+1207603x^3\\
 &\qquad+22827542x^2+202280436x+652159656.
\end{split}
 \tag{F1T28.14}
\]
All denominators in (F1T28.13) are positive.
Thus both alternatives in (F1T28.6) fail under the
assumed row and column bound, proving the first
inequality of (F1T28.1).

T27 gives at least four boxes outside both the first
row and first column. Together with the strict bound,
$4<n/8$, so $n\ge33$. The increasing function $D_3$
then gives $m\ge D_3(33)=4928$.
If a row-based description is wanted, the explicit
sign-conjugate intertwiner already constructed in
T22--T26 sends the actual module to the transposed
Specht model while replacing $c$ by
$c\,\operatorname{sgn}\sigma$. Along this specified
isomorphism the surviving shapes have the form
$(n-k,\nu)$ with $4\le k<n/8$; no independence of
the two original characters is introduced.

\paragraph{Framework and source scope.}
The only new maps are the exact box-to-diagonal-hook
bijection, the explicit content-increasing box moves,
and inequalities on the class-sum eigenvalues of the
original representation. The inertia commutator,
odd-square identity and primitive-corner maps are
the previously checked actual maps of T25.
The source curve, its coefficient local system,
the original sixfold corner, all pull-push scales
and the target Hermitian form remain those of
T24--T27. Consequently its Hilbert/Chow,
Poincare-biextension, cycle-class and Adams/GRR
maps are retained along those identical objects.
This proof makes no replacement by a larger
Jacobian or a whole isotypic factor.

The exact content-character identities are the
Lassalle and class-sum formulas fully checked in
T22/T25; the diagram identities and inequalities
used here are proved above. The finite certificate
checks every partition in degrees22--32 as well
as all symbolic identities and positive-coefficient
expansions. The universal restriction comes from
the displayed proof.
